\par\clearpage\SourceLeaf{302}\ArithmeticHead{217}{R}{}
\HistoricalFoot{X\textsubscript{2} Abh. 2.}{28}
\begin{center}\textbf{Berichtigungen zu der Abh. 2 (Gauss’ Werke Bd. X, 2).}\end{center}
\begingroup
\fontsize{9.2}{13}\selectfont
\setlength{\tabcolsep}{3pt}
\noindent\begin{tabular}{@{}r r p{.81\linewidth}@{}}
Seite & 10, & Zeile 21 lies ihn statt ihnen.\\
» & 13, & Zeile 2 v. u. fehlt Komma vor dass.\\
» & 14, & Zeile 13 v. u. fehlt Komma vor und.\\
» & 25, & Zeile 13 lies ganz statt ganze.\\
» & 31, & Zeile 16 lies vollständigen statt kompletten.\\
» & 38, & Zeile 3 der Fussnote 54, lies im statt in.\\
» & 40, & Zeile 13 fehlt Komma vor was.\\
» & 40, & Zeile 16 ist Komma vor nicht zu streichen.\\
» & 46, & Zeile 1 fehlt Komma vor aus.\\
» & 57, & Zeile 1 lies deren statt dessen.\\
» & 60, & Zeile 22 ist Komma vor \textsc{Gauss} zu streichen.\\
» & 70, & Zeile 14 lies Analogon statt Anaolgon.\\
» & 70, & Zeile 23 ist zwischen (16) und auftretenden einzufügen: mit \(\displaystyle\frac\nu\pi\) multipliziert.\\
» & 72, & Zeile 3 v. u. fehlt Komma vor und.\\
» & 82, & Zeile 4 ist Komma vor nach zu streichen.\\
» & 93, & Zeile 14 lies der statt des.\\
» & 97, & Fussnote 112, Zeile 1 fehlt Komma vor so.\\
» & 124, & Zeile 16 lies unterhalb statt unterhab.\\
» & 140, & Zeile 7 v. u. lies älteren statt ältern.\\
» & 149, & Zeile 8 fehlt Komma vor später, Zeile 2 der Fussnote 216 lies \textit{los} statt \textit{os}.\\
» & 214, & Namenregister lies \textsc{Alembert}, J. statt T.\\
» & 215, & Namenregister lies \textsc{Liouville} statt \textsc{Lionville}.\\
» & 216, & Namenregister lies \textsc{Pfaff}, J. \textsc{Fr.} statt T. \textsc{Fr.}\\
\end{tabular}
\par\bigskip\begin{center}\rule{.23\linewidth}{.3pt}\end{center}
\par\clearpage\endgroup
